\section{Exact rational certificates with objective slack}\label{sec:certificates}

The box hierarchy bounds can supply exact certificates for rational data.
The hierarchy bounds first establish real finite-order membership with a
positive constant margin. An explicit rational interior direction then
allows rounding and exact coefficient correction. We prove a polynomial bound on the encoding
length of the resulting rational Gram matrices in the expanded SDP
dimensions and rational input length, including the slack. This is a
specialization of established strict-feasibility rational
SOS recovery methods, rather than a new general rationalization method.

\subsection{The finite cone and its encoding}\label{sec:rat-encoding}

Fix $r\ge1$ and nonempty continuous bags. Empty bags can have their
constant terms assigned to another bag. Let $\mathcal C_r$ be either
$\sum_b\Pre_r(B_b)$ or $\sum_b\Mod_r(B_b)$. For the preordering take
$\mathcal I_b=\{I\subseteq B_b:\abs I\le r\}$; for the ordinary module take
$\mathcal I_b=\{\varnothing\}\cup\{\{i\}:i\in B_b\}$.
For each retained block define
\[
 q_{b,I}=\prod_{i\in I}(1-x_i^2),\qquad
 v_{b,I}=(x^\beta:\abs\beta\le r-\abs I),\qquad
 s_{b,I}=\binom{v_b+r-\abs I}{v_b}.
\]
These are full monomial bases, including every allowed monomial. Let
\[
 D_{\mathcal C}=\sum_{b,I}s_{b,I},\quad
 V_{\mathcal C}=\sum_{b,I}\frac{s_{b,I}(s_{b,I}+1)}2,\quad
 S_{\mathcal C}=\max_{b,I}s_{b,I}.
\]
Let $M_{\mathcal C}$ be the number of distinct global monomials supported
in some bag and of total degree at most $2r$. Coefficients of coincident
monomials from overlapping bags are collected in the same output row.
For a tuple of symmetric matrices the coefficient map is
\begin{equation}\label{eq:rat-map}
 \mathcal G_{\mathcal C}(Q)=
 \sum_b\sum_{I\in\mathcal I_b}q_{b,I}v_{b,I}^{T}Q_{b,I}v_{b,I}.
\end{equation}
We use its monomial coefficient vector when taking a coefficient norm.
The independent coordinates are the $V_{\mathcal C}$ upper-triangular
entries. Its matrix has integer entries of modulus at most two. Each
column expands into at most $2^w$ monomials; for the ordinary module the
bound is two. Matrices $Q_{b,I}\succeq0$ represent exactly $\mathcal C_r$.

\begin{theorem}[Short exact rational Gram certificates]\label{thm:rational}
Let $p$ have rational coefficients and let $\tau\in\Q_{>0}$. Suppose
\begin{equation}\label{eq:rat-membership}
 p-\tau\in\mathcal C_r
\end{equation}
with real Gram matrices in the full bases above. Then there are rational
symmetric matrices with
\begin{equation}\label{eq:rat-certificate}
 \mathcal G_{\mathcal C}(Q)=p,\qquad
 Q_{b,I}\succeq\frac{\tau}{2D_{\mathcal C}}I
 \quad\text{for every retained block}.
\end{equation}
Their total binary encoding length is polynomial in
$V_{\mathcal C},M_{\mathcal C}$, the rational input length (including
$\tau$ and the bags), and $\log^+(1/\tau)$, where
$\log^+t=\max\{0,\log_2t\}$. The coefficient identity and positive
semidefiniteness of a supplied tuple can be checked by exact rational
arithmetic in polynomial time in its expanded encoding length.
Running intersection is unnecessary for this rationalization statement
once \eqref{eq:rat-membership} is supplied.
\end{theorem}

The real membership prerequisite is stronger than pointwise positivity at
a prescribed order. The theorem gives a polynomial-size witness; rounding
an unspecified real witness in its proof does not assert a polynomial-time
method for locating that witness.

\subsection{An interior certificate of one}\label{sec:rat-interior}

\begin{lemma}[A diagonal rational interior direction]\label{lem:rat-interior}
There is a diagonal rational tuple $H$ in the chosen full bases such that
\begin{equation}\label{eq:rat-interior}
 \mathcal G_{\mathcal C}(H)=1,\qquad
 H_{b,I}\succeq D_{\mathcal C}^{-1}I,\qquad
 \sum_{b,I}\tr H_{b,I}\le r+1.
\end{equation}
Its entries have polynomial encoding length in the expanded dimensions.
\end{lemma}
\begin{proof}
In one bag order the coordinates $1,\dots,v$ and a generator set
$I=\{i_1,\dots,i_k\}$. For each diagonal position indexed by $\beta$,
with $\abs\beta+k\le r$, the following complement identity holds:
\begin{align}
 1-x^{2\beta}q_I
 ={}&\sum_{i=1}^{v}\sum_{j=0}^{\beta_i-1}
 \bigl(x_1^{\beta_1}\cdots x_{i-1}^{\beta_{i-1}}x_i^j\bigr)^2
 (1-x_i^2) \label{eq:rat-complement}\\
 &+\sum_{a=1}^{k}(x^\beta x_{i_a})^2
           \prod_{h<a}(1-x_{i_h}^2).\nonumber
\end{align}
The first telescope is $1-x^{2\beta}$, and the second is
$x^{2\beta}(1-q_I)$. The first sum has total degree at most
$2\abs\beta$. A term in the second uses monomial degree
$\abs\beta+1$ and $a-1$ box generators, so its total degree is
$2(\abs\beta+a)\le2r$. All terms therefore use retained preordering
blocks. If $I$ is empty or singleton, the first sum uses only singleton
blocks and the second uses only an empty-generator block; hence the
identity also stays within the ordinary module.

Write $1=x^{2\beta}q_I+(1-x^{2\beta}q_I)$ for every one of the
$D_{\mathcal C}$ diagonal positions across all bags, sum the displayed
identities, and divide by $D_{\mathcal C}$. The resulting tuple is
rational and diagonal. Each position receives its own baseline
$1/D_{\mathcal C}$; all other contributions are nonnegative. Each source
identity has $1+\abs\beta+\abs I\le r+1$ monomial-square terms counted
with multiplicity, proving the trace bound. Before division the entries
are nonnegative integers at most $D_{\mathcal C}(r+1)$, which proves the
encoding claim.
\end{proof}

\subsection{A rational bound on the feasible Grams}\label{sec:rat-outer}

Let $\ell$ be expectation under the global product uniform probability
measure on the box and put
\[
 W_{b,I}=\ell(q_{b,I}v_{b,I}v_{b,I}^{T}).
\]
The same global measure is used for every bag, so overlaps create no
additional multiplicity in a polynomial identity.

\begin{lemma}[Explicit outer bound]\label{lem:rat-outer}
Write $S=S_{\mathcal C}$ and set
\[
 q_0=((2r+3)!)^{2w},\qquad
 \kappa_0=q_0^{-S}S^{-(S-1)},
\]
with the last factor one for $S=1$. Then $W_{b,I}\succeq\kappa_0 I$.
For $C_p$ the sum of the absolute monomial coefficients of $p$, every
PSD tuple representing $p$ satisfies
\begin{equation}\label{eq:rat-outer}
 \sum_{b,I}\tr Q_{b,I}\le C_p/\kappa_0,
 \qquad \norm Q_F\le C_p/\kappa_0.
\end{equation}
Here $\norm Q_F$ is the aggregate Frobenius norm. The rational number
$R_{\mathrm{out}}=1+C_p/\kappa_0$ has polynomial encoding length in the
expanded dimensions and input length.
\end{lemma}
\begin{proof}
A nonzero polynomial cannot vanish almost everywhere on an open box, and
$q_{b,I}$ is strictly positive there. Thus $W_{b,I}$ is positive definite.
Its entries are products of univariate uniform integrals, which for even
$a$ are
\[
 \ell(x^a)=\frac1{a+1},\qquad
 \ell(x^a(1-x^2))=\frac2{(a+1)(a+3)};
\]
for odd $a$ both are zero. All exponents here are at most $2r$, and $q_0$
clears all denominators. A block of size $s$ consequently has positive
determinant at least $q_0^{-s}$, since $q_0W_{b,I}$ is an integer matrix
with positive integer determinant. Each diagonal is at most one, so
$\tr W_{b,I}\le s$ and every eigenvalue is at most $s$. Dividing its
determinant by the product of the other eigenvalues gives
$\lambda_{\min}(W_{b,I})\ge q_0^{-s}s^{-(s-1)}\ge\kappa_0$.
For PSD matrices, trace pairing with $W_{b,I}\succeq\kappa_0 I$ gives
\[
 \kappa_0\sum_{b,I}\tr Q_{b,I}
 \le\sum_{b,I}\tr(W_{b,I}Q_{b,I})
 =\ell(p)\le C_p.
\]
The aggregate Frobenius norm of a PSD tuple is at most its total trace.
Finally
$\log_2(1/\kappa_0)=O(Swr\log(r+2)+S\log(S+1))$.
A nonempty bag's empty-generator block contains $1,x_i,\ldots,x_i^r$,
so $r+1\le S$; the width is at most the explicit bag-list length.
These observations prove polynomial bit growth in the stated expanded
parameters. The bound can be numerically very conservative.
\end{proof}

\subsection{Exact correction and proof of the theorem}\label{sec:rat-correction}

\begin{proof}[Proof of \cref{thm:rational}]
To simplify the formulas write $D_*=D_{\mathcal C}$,
$V_*=V_{\mathcal C}$ and $\mathcal G=\mathcal G_{\mathcal C}$.
By \eqref{eq:rat-membership} there is a PSD tuple $G$ with
$\mathcal G(G)=p-\tau$. Add the interior direction of
\cref{lem:rat-interior}: $Q^*=G+\tau H$ satisfies
\[
 \mathcal G(Q^*)=p,\qquad Q^*_{b,I}\succeq\varepsilon_* I,
 \qquad\varepsilon_*=\tau/D_*.
\]
\Cref{lem:rat-outer} bounds all its entries by $R_{\mathrm{out}}$.

The coefficient map has a simple rational right inverse. For each global
output monomial $x^\beta$, choose a containing bag and split
$\beta=\alpha+\zeta$ with $\abs\alpha,\abs\zeta\le r$.
Such a split exists because $\abs\beta\le2r$.
In its empty-generator Gram block put one in the diagonal position if
$\alpha=\zeta$, or $1/2$ in the two symmetric off-diagonal positions
otherwise. Call this tuple $Z_\beta$. Then
$\mathcal G(Z_\beta)=x^\beta$ and $\norm{Z_\beta}_F\le1$.
Distinct monomials cannot share a matrix position, since a position has a
unique exponent sum. Thus the choices give simultaneous pivots. Their
linear extension $Z$ satisfies
\begin{equation}\label{eq:rat-right-inverse}
 \mathcal G Z=\operatorname{id},\qquad \norm{Z(c)}_F\le\norm c_1.
\end{equation}
This inverse acts on global collected coefficients, rather than
correcting each bag separately.

Round each upper-triangular entry of $Q^*$ to a multiple of
$h=2^{-B}$ with absolute error at most $h$, preserving symmetry, and call
the tuple $Q^0$. Counting coordinates and generator expansions gives
\[
 \norm{Q^0-Q^*}_F\le2V_*h,\qquad
 \norm{\mathcal G(Q^0-Q^*)}_1\le2^{w+1}V_*h.
\]
Define the rational tuple
\begin{equation}\label{eq:rat-round-correct}
 Q=Q^0+Z(p-\mathcal G(Q^0)).
\end{equation}
Its coefficient identity is exact by \eqref{eq:rat-right-inverse}, and
\[
 \norm{Q-Q^*}_F\le(2+2^{w+1})V_*h\le2^{w+2}V_*h.
\]
Choose an integer $B\ge0$ so that
$2^{-B}\le\varepsilon_*/(2^{w+3}V_*)$.
Then each block perturbation has spectral norm at most
$\varepsilon_*/2$, so every $Q_{b,I}\succeq(\varepsilon_*/2)I$,
which proves \eqref{eq:rat-certificate}.

The initial rounded entries have numerator length
$O(B+\log_2(R_{\mathrm{out}}+1))$. If $L$ is the least common multiple of
$2^B$ and all denominators of coefficients of $p$, then coefficients of
$p-\mathcal G(Q^0)$ have denominators dividing $L$. The correction $Z$
introduces at most a factor of two. Thus $2L$ clears all corrected Gram
entries, including an off-diagonal pivot that halves an input coefficient
whose denominator has a larger power of two than $2^B$. Its encoding
length is at most a constant plus $B$ and the sum of the input denominator
lengths. Coefficient formation uses polynomially many rational additions
and multiplications by integers of modulus at most two in the expanded
map. The correction numerators therefore have polynomial bit length.
Since $B=O(w+\log(V_*+1)+\log(D_*+1)+\log^+(1/\tau))$,
\cref{lem:rat-outer} proves the total witness-size claim.

A supplied tuple is checked by expanding \eqref{eq:rat-map} and comparing
all rational coefficients, then testing each rational symmetric matrix
for PSD by exact elimination. Positive pivots permit Schur complements;
a negative diagonal disproves PSD, and a zero diagonal with a nonzero
off-diagonal entry also disproves PSD. An all-zero residual block is PSD.
The elimination arithmetic has polynomial bit growth by determinant
bounds. Thus both identity matching and PSD testing take polynomial time
in the expanded certificate encoding length.
\end{proof}

Formula \eqref{eq:rat-round-correct} also describes a practical conversion
step for an approximate Gram tuple: round, enforce coefficients exactly,
and check the resulting PSD matrices. For a quantitative sufficient
condition, if an approximate tuple is within aggregate Frobenius distance
$e$ of a tuple $Q^*$ with margin $\varepsilon_*$ and exact identity, its
rounded-and-corrected tuple differs from $Q^*$ by at most
\[
 (1+2^{w+1}\sqrt{V_{\mathcal C}})e
       +2^{w+2}V_{\mathcal C}h.
\]
Indeed the upper-triangular entrywise error sum is at most
$\sqrt{V_{\mathcal C}}e$, and each column of the map has coefficient
one-norm at most $2^{w+1}$. If this displayed bound is at most
$\varepsilon_*/2$, the remaining Gram margin is at least
$\varepsilon_*/2$. An arbitrary numerical SDP output need not satisfy
these conditions. An exact identity and a successful rational PSD check
are a certificate regardless of how the tuple was obtained.

\subsection{From the two hierarchy rates to rational witnesses}
\label{sec:rat-rates}

\begin{corollary}[Rational box certificates at the quantitative orders]
\label{cor:rational-mod}
Let every supplied local objective polynomial $f_b$ have rational
coefficients, and let $\lambda\in\Q$ and $\tau\in\Q_{>0}$.
Choose either of the following cones and rational gap
budgets:
\begin{enumerate}[label=(\alph*)]
\item For $\mathcal C_r=\sum_b\Pre_r(B_b)$, assume
$r\ge\max\{w,d\}$ and set
\[
 E_{\mathrm{pre}}=\frac{3\Abud}{2(\lfloor r/w\rfloor+1)^2+1}.
\]
\item For $\mathcal C_r=\sum_b\Mod_r(B_b)$, use
$s\ge\max\{2,d_\infty\}$, even $N\ge2$, and $r\ge wD$ as in
\cref{thm:mod}. Define the rational majorant
\[
 \eta_{\Q}=\frac{N+1}{C_s}
 \left(\frac{d_\infty^2}{s^2}+\frac{3d_\infty^2}{2s}\right)
       +2(D+1)\delta,
 \qquad E_{\mathrm{mod}}=\Cf\bigl((1+\eta_{\Q})^w-1+2\Delta_w\bigr).
\]
\end{enumerate}
If $\fstar-\lambda\ge E_{\mathcal C}+\tau$, then $f-\lambda$ has the
exact rational Gram certificate of \cref{thm:rational} in that same cone.
The budget is $O(r^{-2})$ in (a) and
$\Cf O_{w,d_\infty}(\log^3r/r^2)$ in (b), for the stated fixed-data
parameter choices. Their encoding length is polynomial in the expanded
finite cone dimensions and rational input length, including the slack.
\end{corollary}
\begin{proof}
All Chebyshev coefficients and budgets are rational for rational monomial
data. In (b), $\sqrt{2(D+1)}\le2(D+1)$ gives $\eta_{\Q}\ge\eta$.
Monotonicity and \cref{thm:mod} show that $E_{\mathrm{mod}}$ is a valid
rational gap bound. In either case \cref{lem:slater} gives real membership
$f-\rho_r\in\mathcal C_r$ with $\rho_r$ the corresponding moment value.
The gap bound and promise imply $\rho_r\ge\lambda+\tau$. Adding this
nonnegative constant to a square block gives
$f-\lambda-\tau\in\mathcal C_r$, precisely the prerequisite of
\cref{thm:rational}.

For (a), the rate is \cref{thm:pre}. For (b), the parameter choice of
\cref{cor:mod-rate} has $D=O_w(s\log s)$ and
$\delta=O_w(s^{-3})$. The replacement term $2(D+1)\delta$ is therefore
$O_w(\log s/s^2)$, the same order as the other terms in $\eta$.
The correction estimates are unchanged, so the rational majorant retains
$O_{w,d_\infty}(\log s/s^2)=O_{w,d_\infty}(\log^3r/r^2)$.
\end{proof}

If the unlabelled problem has no continuous variables, its rational
constant lower-bound certificate can be written directly; no width or
kernel formula is needed. Finite-state certificates of
\cref{prop:finite-dual} are a separate real statement. The rational theorem
above covers box-only unlabelled cones with full bases. It does not extend
automatically to extra polynomial generators, labelled cones, or reduced
Gram bases, which need their own interior directions and bounds.

Peyrl and Parrilo~\cite{peyrl2008-computing-sum-of-squares-decompositions}
already give numerical-symbolic rational SOS recovery by rounding and
projection under strict feasibility. Davis and
Papp~\cite{davis2024-rational-dual-certificates-for-weighted} give quantitative
rational certificate bounds and rational algorithms for weighted SOS,
including discussion of sparse structure. Our contribution here is to
connect the proved sparse hierarchy budgets to real finite-order membership
and to supply elementary box-specific interior, outer, and coefficient
correction bounds. The ordinary-module restriction uses the same proof;
it is not a separate rationalization method. The encoding bound is in
expanded order-dependent Gram dimensions, including all retained generator
products. It gives no polynomial bound in the original optimization input
size and no guarantee about a particular numerical solver's accuracy or
running time.
